% Appendix A — Theorem 1 Boundary Cases
% Analyzes the six mechanisms marked † in Table 1 (§7 design-space map).
% Each entry: (1) mechanism description, (2) why it is a boundary case,
% (3) capability classification with conditions, (4) conditions under which
% classification changes.
%
% REVIEW STATUS: Draft. Espresso and mev-commit entries depend on Task 2.3
% (preconfer codebase evaluation) for accurate technical detail.

\section{Theorem 1 Boundary Cases}
\label{app:boundary}

This appendix analyzes the six mechanisms in Table~\ref{tab:designspace}
marked as boundary cases ($\dagger$). For each mechanism, we state the
conditions under which the classification holds and the conditions under
which it would change.

%-----------------------------------------------------------------------
\subsection{Espresso HotShot Execution Attestation}
\label{app:espresso}

\paragraph{Mechanism.}
Espresso's HotShot consensus provides a shared finality layer for multiple
rollups. In the execution attestation variant, the HotShot committee attests
to the execution outcome of a batch before it is finalized.

\paragraph{Why a boundary case.}
The attestation is capability (i) if and only if it covers the execution
outcomes of \emph{both} rollup STFs jointly. If the attestation covers only
one rollup's execution (or covers each independently without a joint predicate),
then $\phi(B) = \mathsf{true}$ for the attestation does not guarantee that
both $\delta_1$ and $\delta_2$ succeed.

\paragraph{Classification conditions.}
\begin{itemize}
  \item \textbf{Capability (i) holds} if: the HotShot attestation is a joint
    predicate $\phi(\delta_1, \delta_2)$ that is true iff both STFs succeed,
    and the attestation is issued before the commitment $\Gamma(B)$ is
    externally observable.
  \item \textbf{A4 only} if: the attestation covers each rollup independently
    (two separate attestations), or the attestation is issued after $\Gamma(B)$
    is already final (i.e., the attestation is a post-hoc proof, not a
    pre-commitment predicate).
\end{itemize}

\paragraph{Current assessment.}
Based on public Espresso documentation~\cite{espresso-docs}, HotShot provides
finality per rollup batch, not a joint predicate over multi-rollup bundles.
The \emph{shared} component is ordering, not execution outcome. We therefore
classify Espresso as \emph{conditional} capability (i): the classification
would change if a joint execution attestation mechanism were added.

% TODO (Task 2.3): cross-check against current HotShot codebase.

%-----------------------------------------------------------------------
\subsection{Across Protocol (UMA Oracle + Escrow)}
\label{app:across}

\paragraph{Mechanism.}
Across Protocol is an intent-based bridging protocol. A user posts an intent;
a filler executes on the destination chain immediately; the origin-chain funds
are released from escrow after UMA oracle confirmation. The UMA Optimistic
Oracle validates the fill before releasing origin funds.

\paragraph{Why a boundary case.}
Across uses UMA, an optimistic oracle with a challenge period. During the
challenge window, the filler's origin-chain payment is not yet final — the
UMA oracle can dispute the fill. This is capability (ii): the ``irreversible
effect'' on the origin chain (fund release) is prevented until confirmation.
However, the destination-chain fill is immediately final (no escrow on the
destination side). This asymmetry makes the classification non-trivial.

\paragraph{Classification conditions.}
\begin{itemize}
  \item \textbf{Capability (ii) holds for the origin-side effect}: The fund
    release on the origin chain is blocked until UMA confirmation. If the fill
    is disputed, origin funds are not released — partial-effect prevention holds
    for the origin leg.
  \item \textbf{Destination leg is not protected}: The filler's execution on
    the destination chain is immediately final. If UMA disputes the fill after
    destination execution, the filler bears the loss (an A4 mechanism, not A2).
  \item \textbf{Full A2 for users}: From the \emph{user's} perspective, funds
    are either moved completely (both legs final) or the origin escrow is
    returned. This is A2 for user funds.
  \item \textbf{Not A2 for fillers}: The filler's capital is exposed during the
    challenge window. A filler can be penalized even after executing correctly.
    This is an A4 risk for the filler.
\end{itemize}

\paragraph{Classification.}
Across provides capability (ii) for user-facing fund atomicity, with a residual
A4 exposure for fillers during the UMA challenge window. In the context of
\S\ref{sec:designspace}, we classify Across as capability (ii) for the bundle
as a whole, noting the filler-side A4 residual. This is the correct
classification for determining whether a PEFO attacker could exploit the
application: the user's funds are A2-protected, so the application is not a
valid PEFO target for user-fund extraction.

%-----------------------------------------------------------------------
\subsection{UniswapX Cross-Chain Signed Orders}
\label{app:uniswapx}

\paragraph{Mechanism.}
UniswapX uses signed order structures where a user signs an order specifying
input token, output token, and a deadline. A filler executes the order
on-chain. For cross-chain variants, the order may span two chains.

\paragraph{Why a boundary case.}
The signed-order constraint provides a form of capability (ii): the filler
cannot execute a partial order (they must provide the full output amount or
the transaction reverts). However, the cross-chain synchronization still
depends on filler economics (they fill on chain A, then claim on chain B),
which is A4.

\paragraph{Classification.}
Capability (ii) for the fill atomicity guarantee (within a single chain, the
order is atomic). A4 for the cross-chain claim timing. Combined classification:
(ii)+A4. The cross-chain synchronization gap is covered by filler incentives,
not by a protocol-level capability (i) or (ii) spanning both chains.

%-----------------------------------------------------------------------
\subsection{Shutter Network (Threshold Encryption)}
\label{app:shutter}

\paragraph{Mechanism.}
Shutter uses threshold encryption: transactions are encrypted before being
submitted to the sequencer. The decryption key is released only after the
block is committed. This prevents frontrunning and selective exercise.

\paragraph{Why a boundary case.}
Shutter provides both (i) and (ii) simultaneously, which is stronger than
Theorem~\ref{thm:a2-necessary} requires. The hybrid classification warrants
explanation.

\paragraph{Capability (i) analysis.}
The sequencer commits to the \emph{encrypted} bundle $\text{Enc}(B)$ before
seeing the plaintext transactions. This prevents the sequencer from selectively
including or excluding transactions based on execution outcomes. Once the
decryption key is released and the transactions are decrypted, both $\tau_1$
and $\tau_2$ execute atomically in the committed ordering. The pre-execution
predicate is: ``both transactions will execute as committed'' — which holds by
construction because the sequencer cannot modify the committed order after
decryption.

\paragraph{Capability (ii) analysis.}
The threshold decryption requirement prevents partial execution: if the
threshold committee withholds decryption (e.g., due to liveness failure), both
legs are delayed together — not selectively. An attacker cannot cause $\tau_1$
to execute while preventing $\tau_2$.

\paragraph{Liveness caveat.}
Shutter's A2 guarantee is contingent on threshold committee liveness. If the
committee fails to release the decryption key within the timeout, both legs
fail together (preserving A2) but the protocol halts (a liveness failure). This
is a standard property of threshold cryptography and does not affect the safety
classification.

\paragraph{Classification.} (i)+(ii), conditional on threshold committee liveness.

%-----------------------------------------------------------------------
\subsection{EigenLayer Restaking (AVS Preconfirmation)}
\label{app:eigenlayer}

\paragraph{Mechanism.}
EigenLayer allows operators to restake ETH and use it as collateral for Active
Validator Services (AVS). An AVS providing preconfirmations can slash restaked
operators who violate commitment terms.

\paragraph{Why a boundary case.}
Restaking extends the bond mechanism to larger stake. The question is whether
the larger stake creates a new capability (i) or (ii), or whether it is
fundamentally still A4.

\paragraph{Analysis.}
Restaking increases the size of $B_{\max}$ in Theorem~\ref{thm:bond-insufficient}.
It does not:
\begin{itemize}
  \item Provide a success predicate $\phi(B)$ (the operator cannot evaluate
    $\delta_1, \delta_2$ before the commitment — they are still lazy by
    Definition~\ref{def:lazy});
  \item Provide revert authority over either rollup's state (slashing removes
    stake from the operator, it does not revert execution on $R_1$ or $R_2$).
\end{itemize}
Therefore, restaking is A4 only. The bound in Theorem~\ref{thm:bond-insufficient}
applies: for any finite $B_{\max}$ (however large), there exists an option value
$V > B_{\max}$ with nonzero probability, making violation profitable.

\paragraph{Classification.} A4 only, regardless of stake size.

%-----------------------------------------------------------------------
\subsection{mev-commit / Bolt Preconfirmation}
\label{app:mevcommit}

\paragraph{Mechanism.}
mev-commit (Primev) and Bolt are preconfirmation protocols where validators
(or searchers acting as preconf providers) issue signed commitments to include
specific transactions in a future block. The commitment is backed by a bond;
failure to include results in slashing.

\paragraph{Why a boundary case.}
mev-commit and Bolt provide stronger guarantees than a simple bond: the signed
preconfirmation is a named commitment to a specific transaction, verifiable
on-chain. The question is whether the named commitment creates capability (i)
or (ii).

\paragraph{Analysis.}
A mev-commit preconfirmation commits to \emph{inclusion} of a transaction, not
to the \emph{execution outcome}. The preconf provider signs ``transaction $\tau$
will be included in block $B$''; they do not sign ``$\delta(\sigma, \tau)$ will
succeed.'' Specifically:
\begin{itemize}
  \item \textbf{No capability (i)}: The preconf provider does not evaluate the
    STF before issuing the commitment. The signed preconf is not a success
    predicate $\phi(\tau)$.
  \item \textbf{No capability (ii)}: The preconf provider does not have authority
    to revert the rollup state if execution fails after inclusion. Slashing
    punishes the preconf provider, but $\pi_1(\tau_1)$ (the irreversibility point)
    is determined by the rollup protocol, not the preconf provider.
\end{itemize}

\paragraph{Bond-slashing model.}
The mev-commit bond is slashed if the committed transaction is not included.
This is a compensation mechanism for inclusion failure, not an execution
atomicity guarantee.

% TODO (Task 2.3): cross-check against current mev-commit and Bolt codebase
% for any features added after documentation was written.

\paragraph{Classification.} A4 only.
